\chapter{Anhedonia (Loss of Interest)}

\section*{Definition}
Anhedonia is a reduced ability to experience pleasure, interest, or motivation in activities that were previously rewarding. It is common in depression but can also occur with schizophrenia, trauma-related disorders, substance use, neurologic disease, medication effects, chronic illness, or grief. It is a symptom, not a diagnosis.

\section*{When to See a Doctor}

\begin{itemize}
 \item Loss of pleasure or interest lasting two weeks or more, interfering with work, relationships, self-care, or sleep
 \item Any suicidal thoughts, self-harm urges, psychosis, inability to care for basic needs, or rapidly worsening symptoms -- seek urgent help
\end{itemize}

\section*{How It Develops}
Anhedonia reflects changes in reward processing, motivation, learning, and emotional experience. Brain circuits involving the striatum, prefrontal regions, and other networks may be involved, but no single neurotransmitter or brain scan explains every case. Sleep, stress, illness, medications, substances, and hormonal changes can all influence these systems.

\section*{Causes}
\begin{itemize}
 \item Major depressive disorder (most common association)
 \item Perimenopausal and menopausal hormonal changes
 \item Chronic stress, trauma, grief, or prolonged sleep disruption
 \item Schizophrenia or schizoaffective disorder
 \item Substance use disorders (dopamine receptor downregulation)
 \item Parkinson disease (dopamine depletion)
 \item Post-traumatic stress disorder
 \item Medication side effects (SSRIs, antipsychotics)
 \item Hypothyroidism
 \item Medical illness, chronic pain, or medication effects; thyroid disease, anemia, or nutritional deficiency may mimic or worsen depressive symptoms
\end{itemize}

\section*{Related Symptoms}
\begin{itemize}
 \item Low mood, guilt, hopelessness, irritability, or anxiety
 \item Fatigue, sleep or appetite changes, and difficulty concentrating
 \item Social withdrawal, reduced motivation, or loss of emotional responsiveness
 \item Suicidal thoughts or thoughts of death
\end{itemize}


\section*{Medical Evaluation}
\begin{itemize}
 \item Your doctor will take a detailed psychiatric history including duration of anhedonia, associated mood symptoms, sleep and appetite changes, and suicidal ideation.
 \item Screening questionnaires such as the PHQ-9 can assess depressive symptoms; the Snaith-Hamilton Pleasure Scale may help quantify anhedonia but does not diagnose its cause.
 \item Tests are guided by the history and examination. A clinician may check a blood count, thyroid function, B12, glucose, or other studies when anemia, endocrine disease, medication effects, or another medical cause is possible; routine vitamin D testing is not required for everyone.
 \item Referral to a psychiatrist or clinical psychologist is arranged for comprehensive assessment and psychotherapy, particularly if depression or another mood disorder is suspected.
\end{itemize}

\section*{Treatment Options}
\begin{itemize}
 \item Treat the underlying condition. For depression, psychological therapy, behavioural activation, antidepressants, or a combination may be offered according to severity, preference, and clinical history.
 \item If symptoms persist, a clinician may adjust treatment or consider options such as bupropion or antipsychotic augmentation in selected cases; these medicines have important risks and are not treatments for anhedonia alone.
 \item Behavioural activation therapy, a structured CBT approach that schedules rewarding activities, directly targets anhedonia by rebuilding reward circuitry.
 \item During perimenopause, hormone therapy may be considered for depressive symptoms that begin with other menopause symptoms and do not meet criteria for a depressive disorder; treatment requires an individualized risk assessment.
\end{itemize}

\section*{Self-Care and Home Management}
\begin{itemize}
 \item Establish a daily routine with small, achievable goals such as a short walk, a hobby for 10 minutes, or completing one household task.
 \item Prioritise sleep hygiene: maintain consistent bedtimes, limit screen exposure before sleep, and avoid alcohol or caffeine in the evening.
 \item Engage in social connection even when it feels difficult — brief conversations with trusted friends or family can reduce isolation.
 \item Postpone major decisions if concentration, judgment, or safety is impaired, and discuss concerns with a trusted person or clinician.
 \item Track mood and interest levels daily to identify patterns and share these with your healthcare provider.
\end{itemize}

\section*{References}
\begin{thebibliography}{99}
\bibitem{anhedonia:ref1} Der-Avakian A, Markou A. ``The neurobiology of anhedonia and other reward-related deficits.'' \textit{Trends Neurosci}. 2012;35(1):68-77.
\bibitem{anhedonia:ref2} Treadway MT, Zald DH. ``Reconsidering anhedonia in depression.'' \textit{Neurosci Biobehav Rev}. 2011;35(3):537-555.
\bibitem{anhedonia:ref3} Snaith RP, Hamilton M, et al. ``A scale for the assessment of hedonic tone.'' \textit{Br J Psychiatry}. 1995;167(1):99-103.
\bibitem{anhedonia:ref4} Sadock BJ, Sadock VA, et al. \textit{Kaplan \& Sadock\'s Comprehensive Textbook of Psychiatry}, 10th ed. Wolters Kluwer, 2017.
\bibitem{anhedonia:ref5} UpToDate. ``Approach to the patient with anhedonia.'' Wolters Kluwer, 2023.
\bibitem{anhedonia:ref6} National Institute for Health and Care Excellence. Depression in adults: treatment and management (NG222). Updated 2022.
\bibitem{anhedonia:ref7} National Institute of Mental Health. Depression. Updated 2025.
\end{thebibliography}
